\documentclass[10pt,a4paper]{amsart}
\usepackage[margin=19mm]{geometry}
\usepackage{amsmath,amssymb,mathtools}
\usepackage[scr=rsfs,cal=euler]{mathalfa}
\usepackage{xfrac}
\usepackage[unicode,hidelinks,bookmarks=false]{hyperref}
\newcommand{\ie}{i.e.\ }
\newcommand{\cf}{cf.\ }
\newcommand{\resp}{respectively\ }
\newcommand{\Subsection}{\subsection}
\providecommand{\nobreakdash}{\nobreak\hbox{-}}
\setlength{\emergencystretch}{2em}
\multlinegap0pt
\usepackage{bm}
\usepackage{bbm}
\usepackage{dsfont}
\usepackage{xspace}
\usepackage{cancel}
\usepackage{mathtools}
\usepackage[shortlabels]{enumitem}
\usepackage{amsthm}
\makeatletter
\@ifundefined{theorem}{\newtheorem{theorem}{Theorem}}{}
\@ifundefined{proposition}{\newtheorem{proposition}[theorem]{Proposition}}{}
\makeatother
\title[$\mathrm {U}\_{q\tilde q}\mathfrak{sl}(2;\mathbb R)$ Turaev-V]{$\mathrm {U}\_{q\tilde q}\mathfrak{sl}(2;\mathbb R)$ Turaev-Viro invariants for cusped $3$-manifolds}
\author{Tianyue Liu, Shuang Ming, Xin Sun, Baojun Wu,, Tian Yang}
\date{Source: arXiv:2608.16560v1 (CC BY 4.0)}
\begin{document}
\maketitle

\noindent\emph{Source and licence.} $\mathrm {U}_{q\tilde q}\mathfrak{sl}(2;\mathbb R)$ Turaev-Viro invariants for cusped $3$-manifolds. Version arXiv:2608.16560v1 (CC BY 4.0), distributed under the Creative Commons Attribution 4.0 International licence, as recorded in the arXiv licence metadata for this exact version and shown on the abstract page https://arxiv.org/abs/2608.16560v1. Full rights notice: licenses/arxiv-2608.16560-CC-BY-4.0.txt.\par\smallskip

\noindent\textit{Editorial scope.} Counted unit: the Proposition whose source label is \texttt{critical2} in the article (source0.tex lines 1043--1051, source label \texttt{critical2}), delivered together with the article's own complete original proof of it (source0.tex lines 1053--1147, one \texttt{proof} environment of 95 lines). No mathematics is written, repaired or re-derived: statement and proof are the article's own text line for line. Required context retained uncounted: co-sch (lines 425--427); det- (lines 445--447); cos (lines 450--452). Every definition, display and label of those lines is retained unchanged. Omitted from this package and not redistributed: 1--424, 428--444, 448--449, 453--1042, 1052--1052, 1148--2134. Nothing inside a retained excerpt is omitted or altered: every retained body is delivered exactly as the article's lines are written, so no omission receipt of any kind accompanies this package. Citations: the retained excerpts carry no citation command at all, so no bibliography entry of the article is transcribed and no reference is redistributed. Reference adaptation: none was needed and none was made; every reference inside the delivered unit resolves inside it, so no unresolved reference or citation is produced. Printed slips kept literal: no printed word, symbol, hypothesis or number of the counted statement or of its proof was altered. Layout and class: the article's own class is \texttt{article} with packages including amssymb, amsmath, bm, amsthm, xcolor, textcomp, enumerate, graphicx, caption, subcaption, url, tabu, adjustbox, mathrsfs; the wrapper is the lane's pinned \texttt{amsart} editorial wrapper at 10pt on A4, which restates the theorem-like environments the retained text needs and adds the article's own macro definitions that the retained text needs (none), with extra wrapper packages (bm, bbm, dsfont, xspace, cancel, mathtools, enumitem). The delivered numbering is the unit's own. Assets: no retained span contains a figure environment, a graphics-inclusion command, a PDF-page inclusion, a table or a TikZ picture; no image file is copied into this package, the package declares no asset dependency and no image file is redistributed with it. Licence: version arXiv:2608.16560v1 of this article is distributed under the Creative Commons Attribution 4.0 International licence, as recorded in the arXiv licence metadata for this exact version and shown on the abstract page https://arxiv.org/abs/2608.16560v1. A full rights notice is delivered at licenses/arxiv-2608.16560-CC-BY-4.0.txt. Characters: every retained excerpt is pure ASCII, so the delivered engine, XeLaTeX with its default Latin Modern text font, reports no missing character.
\begin{equation}\label{co-sch}
    \frac{\partial \mathrm{Cov}(\Delta)}{\partial x_{ij}}=\frac{\theta_{ij}}{2}
\end{equation}

\begin{equation}\label{det-}
   \det\mathrm{Gram}(\Delta)<0.
\end{equation}

\begin{equation}\label{cos}
    \cos\theta_{ij}=\frac{G_{kl}}{\sqrt{G_{kk}G_{ll}}},
\end{equation}

\begin{proposition}\label{critical2}
    If $\boldsymbol x=(x_1,\dots, x_6)\in\mathbb R^6$  are the edge lengths of a decorated ideal hyperbolic tetrahedron $\Delta,$ then the function $V_{\boldsymbol x}(\zeta)$ has a unique critical point $\zeta^*$ in the domain $D$ with 
    \begin{equation}\label{12}
    \mathrm{Re}\zeta^*\leqslant \frac{\pi}{12}   
    \end{equation}
    and with critical value
    $$V_{\boldsymbol x}(\zeta^*)=-2\mathbf i \mathrm{Cov}(\Delta),$$
    where $\mathrm{Cov}$ is the co-volume of $\Delta.$
\end{proposition}

\begin{proof} We first show that $V_{\boldsymbol x}$ has a unique  critical point in $D.$  For $i\in\{1,2,3\},$ let $u_i=e^{2y_i},$ and let $z=e^{2\mathbf i\zeta}.$ Then by a direct computation, we have
\begin{equation}\label{eq:mod4}
\frac{d V_{\boldsymbol x}}{d \zeta}=2\mathbf i\log\frac{(1-zu_1)(1-zu_2)(1-zu_3)}{-z^3u_1u_2u_3}\quad\quad (\mathrm{mod}\ 4\pi).
\end{equation} 
Here the $\mathrm{mod}\ 4\pi$ ubiquity  comes from the  $2\mathbf i\log.$ Then, as a necessary condition, $\frac{d V_{\boldsymbol x}}{d \zeta}=0$ implies that 
$\frac{(1-zu_1)(1-zu_2)(1-zu_3)}{-z^3u_1u_2u_3}=1,$ which is equivalent to the following quadratic equation 
$$Az^2+Bz+C=0$$
with $A=u_1u_2+u_1u_3+u_2u_3,$ $B=-(u_1+u_2+u_3),$ and $C=1.$
By a direct computation and (\ref{det-}), we have 
\begin{equation}\label{detGram}
    B^2-4AC = 16\det\mathrm{Gram}(\boldsymbol x)<0.
\end{equation}
Let 
\begin{equation}\label{xi}
z^*=e^{2\mathbf i\zeta^*}=\frac{-B+\sqrt{B^2-4AC}}{2A}\quad\text{and}\quad {z^{**}}=e^{2\mathbf i{\zeta^{**}}}=\frac{-B-\sqrt{B^2-4AC}}{2A}
\end{equation}
be the two roots of this quadratic equation. Then by~\eqref{eq:mod4} we have 
\begin{equation}\label{kk'}
    \frac{d V_{\boldsymbol x}}{d \zeta}\Big|_{\zeta=\zeta^*}=4k\pi\quad\text{and}\quad\frac{d V_{\boldsymbol x}}{d \zeta}\Big|_{\zeta={\zeta^{**}}}=4k'\pi
\end{equation}
for some integers $k$ and $k'.$ We claim that, both $\zeta^*$ and $\zeta^{**}$ are  critical points of $V_{\boldsymbol x},$ among the two only $\zeta^*$ is in $D.$   Indeed, at the special case $x_1=\dots=x_6=0,$ ie., $u_1=u_2=u_3=1,$ we have 
$A=3$ and $B=-3.$ Then $z^*=\frac{3+\sqrt{3}\mathbf i}{6}$ and ${z^{**}}=\frac{3-\sqrt{3}\mathbf i}{6}.$ Hence we can choose $\zeta^*=\frac{1}{2\mathbf i}\log \frac{3+\sqrt{3}\mathbf i}{6}\in D;$ and  there is no ${\zeta^{**}}\in D$ making $e^{2\mathbf i{\zeta^{**}}}=\frac{3-\sqrt{3}\mathbf i}{6},$ as $\mathrm{Im}e^{2\mathbf i\zeta}>0$ for all $\zeta\in D.$ A direct computation also shows that at this special case $k=0,$ i.e.,
$$\frac{d V_{\boldsymbol x}}{d\zeta}\Big|_{\zeta=\frac{1}{2\mathbf i}\log \frac{3+\sqrt{3}\mathbf i}{6}}=0.$$ 
Let $\zeta^{**}=\frac{1}{2\mathbf i}\log \frac{3-\sqrt{3}\mathbf i}{6}.$ Then a direct computation shows that  in this special cases $k'=0$ also, i.e.,
$$\frac{d V_{\boldsymbol x}}{d\zeta}\Big|_{\zeta=\frac{1}{2\mathbf i}\log \frac{3-\sqrt{3}\mathbf i}{6}}=0.$$ 
Now for a general $\boldsymbol x,$  we have $A>0,$ $B<0$ and $B^2-4AC<0;$ and by a direct computation, we have by Weitzenb\"ock's inequality that  $\sqrt{4AC-B^2}=\sqrt{-u_1^2-u_2^2-u_3^2+2u_1u_2+2u_1u_3+2u_2u_3}\leqslant \frac{1}{\sqrt{3}}(u_1+u_2+u_3)=\frac{-B}{\sqrt 3}.$ As a consequence, we have $0<\mathrm{Im}z^*\leqslant \frac{\mathrm{Re}z^*}{\sqrt{3}}.$ Hence $\arg z^* \in (0,\frac{\pi}{6}],$ and  we can choose a unique $\zeta^*$ in $D$ with $\mathrm{Re}\zeta^*\leqslant \frac{\pi}{12}.$  For this choice of $\zeta^*,$ we let $\zeta^{**}=-\overline{\zeta^*}$ so that $e^{2\mathbf i\zeta^{**}}=z^{**}$ and $\mathrm{Re}\zeta^{**}\in (-\frac{\pi}{6},0).$
Since $k$ and $k'$ are integer valued continuous functions on the connected space of the edge lengths of decorated ideal hyperbolic tetrahedra, they are constantly $0,$ ie., 
\begin{equation}\label{=0}
\frac{d V_{\boldsymbol x}}{d \zeta}\Big|_{\zeta=\zeta^*}=\frac{d V_{\boldsymbol x}}{d \zeta}\Big|_{\zeta=\zeta^{**}}=0,
\end{equation}
and $\zeta^*$ is the unique critical point of $V_{\boldsymbol x}$ in  $D.$
\bigskip

Next, we compute the value of  $V_{\boldsymbol x}(\zeta^*).$ To this end, for $\boldsymbol x=(x_1,\dots,x_6)\in \mathbb R^6$ and $\zeta\in D,$ we let 
$V(\boldsymbol x,\zeta)\doteq V_{\boldsymbol x}(\zeta),$ and define 
\begin{equation}\label{W}
W(\boldsymbol x)=V({\boldsymbol x},\zeta^*).
\end{equation}

We first verify that 
\begin{equation}\label{co-sch2}
    \frac{\partial W(\boldsymbol x)}{\partial x_k}=-\mathbf i\theta_k
\end{equation}
for each $k\in\{1,\dots,6\};$ and by the symmetry of the indices, it suffices to consider the case that $k=1.$ 
We claim that 
\begin{equation}\label{pW}
\frac{\partial W(\boldsymbol x)}{\partial x_1}=\frac{1}{2}\log\frac{G_{34}-\sqrt{G_{34}^2-G_{33}G_{44}}}{G_{34}+\sqrt{G_{34}^2-G_{33}G_{44}}},
\end{equation}
where $G_{ij}$ is the $ij$-th cofactor of the Gram matrix $\mathrm{Gram}(\Delta).$ Then by (\ref{cos}), we have 
$$\frac{\partial W(\boldsymbol x)}{\partial x_1}=\frac{1}{2}\log e^{-2\mathbf i\theta_1}=-\mathbf i\theta_1.$$
To prove (\ref{pW}), we let $\zeta^{**}=-\overline {\zeta^*}$ so that $z^{**}=e^{2\mathbf i\zeta^{**}},$ and we define 
$$U(\boldsymbol x)= V(\boldsymbol x,\zeta^{**}).$$
We  will prove that 
\begin{enumerate}[(1)]
    \item  
    $$\frac{\partial (W(\boldsymbol x)-U(\boldsymbol x))}{\partial x_1}=\log\frac{G_{34}-\sqrt{G_{34}^2-G_{33}G_{44}}}{G_{34}+\sqrt{G_{34}^2-G_{33}G_{44}}},$$
    and 

    \item   $$\frac{\partial (W(\boldsymbol x)+U(\boldsymbol x))}{\partial x_1}=0,$$
\end{enumerate}
from which (\ref{pW}) follows immediately. To see (1) and (2), considering both $\zeta^*$ and $\zeta^{**}$ as functions of $\boldsymbol x,$ and by the Chain Rule and  (\ref{=0}), we have
$$\frac{\partial W(\boldsymbol x)}{\partial x_1}=\frac{\partial V}{\partial x_1}\bigg|_{\zeta=\zeta^*}+\frac{\partial V}{\partial \zeta}\bigg|_{\zeta=\zeta^*}\cdot\frac{\partial \zeta^*}{\partial x_1}=\frac{\partial V}{\partial x_1}\bigg|_{\zeta=\zeta^*}+\frac{d V_{\boldsymbol x}}{d \zeta}\bigg|_{\zeta=\zeta^*}\cdot\frac{\partial \zeta^*}{\partial x_1}=\frac{\partial V}{\partial x_1}\bigg|_{\zeta=\zeta^*}$$
and
$$\frac{\partial U(\boldsymbol x)}{\partial x_1}=\frac{\partial V}{\partial x_1}\bigg|_{\zeta=\zeta^{**}}+\frac{\partial V}{\partial \zeta}\bigg|_{\zeta=\zeta^{**}}\cdot \frac{\partial \zeta^{**}}{\partial x_1}=\frac{\partial V}{\partial x_1}\bigg|_{\zeta=\zeta^{**}}+\frac{d V_{\boldsymbol x}}{d \zeta}\bigg|_{\zeta=\zeta^{**}}\cdot\frac{\partial \zeta^{**}}{\partial x_1}=\frac{\partial V}{\partial x_1}\bigg|_{\zeta=\zeta^{**}}.$$
Then we have 
\begin{equation*}
\begin{split}
\frac{\partial (W(\boldsymbol x)-U(\boldsymbol x))}{\partial x_1} = & \frac{\partial V}{\partial x_1}\bigg|_{\zeta=\zeta^*}-\frac{\partial V}{\partial x_1}\bigg|_{\zeta=\zeta^{**}}\\ 
=& \log \frac{z^{**}(1-z^*u_1)}{z^*(1-z^{**}u_1)}\quad\quad (\mathrm{mod}\ 2\pi\mathbf i)\\
=& \log \frac{-u_1+u_2+u_3 - \sqrt{u_1^2+u_2^2+u_3^2-2u_1u_2-2u_1u_3-2u_2u_3}}{-u_1+u_2+u_3 +\sqrt{u_1^2+u_2^2+u_3^2-2u_1u_2-2u_1u_3-2u_2u_3}}\\
=&  \log\frac{G_{34}-\sqrt{G_{34}^2-G_{33}G_{44}}}{G_{34}+\sqrt{G_{34}^2-G_{33}G_{44}}} \quad\quad (\mathrm{mod}\ 2\pi\mathbf i),
\end{split}
\end{equation*}
and 
\begin{equation*}
\begin{split}
\frac{\partial (W(\boldsymbol x)+U(\boldsymbol x))}{\partial x_1} = & \frac{\partial V}{\partial x_1}\bigg|_{\zeta=\zeta^*}+\frac{\partial V}{\partial x_1}\bigg|_{\zeta=\zeta^{**}}\\ 
=& \log \frac{(1-z^*u_1)(1-z^{**}u_1)}{z^*z^{**}u_2u_3}\quad\quad (\mathrm{mod}\ 2\pi\mathbf i)\\
=&\, 0 \quad\quad (\mathrm{mod}\ 2\pi\mathbf i),
\end{split}
\end{equation*}
where the last equality comes from that 
$z^*+z^{**}=-\frac{B}{A}=\frac{u_1+u_2+u_3}{u_1u_2+u_1u_3+u_2u_3}$ and $z^*z^{**}=\frac{C}{A}=\frac{1}{u_1u_2+u_1u_3+u_2u_3}.$ To take care of the  $\mathrm{mod}\ 2\pi\mathbf i$ ambiguity, a direct computation at $\boldsymbol x=(0,\dots,0)$ shows that 
$$\frac{\partial (W(\boldsymbol x)-U(\boldsymbol x))}{\partial x_1} =\log\frac{G_{34}-\sqrt{G_{34}^2-G_{33}G_{44}}}{G_{34}+\sqrt{G_{34}^2-G_{33}G_{44}}}=-\frac{2\pi \mathbf i}{3},$$
and 
$$\frac{\partial (W(\boldsymbol x)+U(\boldsymbol x))}{\partial x_1} = 0.$$
This completes the proof of (1) and (2).



Now, by (\ref{co-sch}) and (\ref{co-sch2}), we have $W(\boldsymbol x)=-2\mathbf i \mathrm{Cov}(\Delta)+K$ for some constant $K.$ By a direct computation, 
$$W(0,\dots,0)= -3\mathbf i\mathrm {D}\bigg(\frac{3+\sqrt{3}\mathbf i}{6}\bigg)=-2\mathbf i\mathrm D\bigg(\frac{1+\sqrt{3}\mathbf i}{2}\bigg)=-2\mathbf i\mathrm{Cov}(0,\dots,0),$$
where $\mathrm D$ is the Bloch-Wigner dilogarithm function defined by $\mathrm D(z)=\mathrm{Im}\mathrm{Li}_2(z)+\arg(1-z)\log|z|$ that computes the volume of the ideal hyperbolic tetrahedron of shape parameter $z,$ and the second equality comes from that $\frac{3+\sqrt{3}\mathbf i}{6}$ is the shape parameter of the ideal hyperbolic tetrahedron with dihedral angles $\big(\frac{\pi}{6},\frac{\pi}{6},\frac{2\pi}{3}\big),$ $\frac{1+\sqrt{3}\mathbf i}{2}$ is the shape parameter of the regular ideal hyperbolic tetrahedron, and the two sides differ by a geometric $3$-$2$ Pachner Move. This shows that the constant $K=0,$ and hence 
$$V_{\boldsymbol x}(\zeta^*)=W(\boldsymbol x)=-2\mathbf i \mathrm{Cov}(\Delta).$$   
\end{proof}

\end{document}
